\documentclass[11pt,a4paper,twoside]{tesis}
% SI NO PENSAS IMPRIMIRLO EN FORMATO LIBRO PODES USAR
%\documentclass[11pt,a4paper]{tesis}

\usepackage{graphicx}
\usepackage[utf8]{inputenc}
\usepackage[spanish]{babel}
\usepackage[left=3cm,right=3cm,bottom=3.5cm,top=3.5cm]{geometry}

\begin{document}

%%%% CARATULA

\def\autor{Pablo Agustín Artuso \\ LU: 282/11 \\ artusopablo@gmail.com}
\def\tituloTesis{Analyzing Windows Telemetry component}
\def\runtitulo{Analyzing Windows Telemetry component}
\def\runtitle{}
\def\director{Rodolfo Baader $<$rbaader@dc.uba.ar$>$}
\def\codirector{Aleksandar Milenkoski $<$amilenkoski@ernw.de$>$}
\def\lugar{Buenos Aires, 2018}
\newcommand{\HRule}{\rule{\linewidth}{0.2mm}}
%
\thispagestyle{empty}

\begin{center}\leavevmode

\vspace{-2cm}

\begin{tabular}{l}
\includegraphics[width=2.6cm]{images/logofcen.pdf}
\end{tabular}


{\large \sc Universidad de Buenos Aires

Facultad de Ciencias Exactas y Naturales

Departamento de Computaci\'on}

\vspace{6.0cm}

%\vspace{3.0cm}
%{
%\Large \color{red}
%\begin{tabular}{|p{2cm}cp{2cm}|}
%\hline
%& Pre-Final Version: \today &\\
%\hline
%\end{tabular}
%}
%\vspace{2.5cm}

\begin{huge}
\textbf{\tituloTesis}
\end{huge}

\vspace{2cm}

{\large Tesis de Licenciatura en Ciencias de la Computaci\'on}

\vspace{2cm}

{\Large \autor}

\end{center}

\vfill

{\large

{Director: \director}

\vspace{.2cm}

{Codirector: \codirector}

\vspace{.2cm}

\lugar
}

\newpage\thispagestyle{empty}


\section*{Abstract (english version)}
\section*{Abstract (spanish version)}
\section*{Greetings}
\section*{Motivation}

\section*{Table of content}

\section*{Introduction}
Explain that this study was carried out in a particular version of windows (1607) 64 bits Enterprise, which lacked of different things (for instance: documentation for different events)
\subsection*{Basic concepts}
\subsubsection*{Reverse Engineering}
Key concepts about what RE means, from a general perspective, how it should tackled which tools are usually there.  64 bits, calling convention .. 
\subsubsection*{Debugging}
Explaniation of the concept of debugging, what is the different with doing static RE. Some words specifically for KERNEL debugging. 
\subsection*{Tools}
\subsubsection*{IDA pro}
Introduction to IDA pro. Explaniation of what it is and how it works.
\subsubsection*{WinDBG}
Introduction to WINDBG pro. Explaniation of what it is and how it works.
\subsubsection*{YARA ? }
\subsubsection*{XPERF? }

\subsection*{Windows components}
\subsubsection*{Event Tracing for Windows}
Complete explainiation of how it works due to its importance for the rest of the analysis. Different components: session, providers, consumers ,etc .
\subsubsection*{Telemetry}
Full description of the different features / characteristic which are involved in this analysis. Explaniation of how the worflow of the data is followed.


\section*{Previous Work}
Some lines about previous works in this topic. Most of them focused on just analysis from traffic / documentation. 

\section*{Analysis}
Specifying the goals we want to achieve and how to get them. 
\subsection*{Understanding how Telemtry makes use of ETW}
How ETW works internally: Functions that are being called, 
\subsubsection*{When and how providers are registered}
\subsubsection*{How writes are carried out}
\subsubsection*{Relation between ETW session and ETW providers}
\subsubsection*{Identifying the buffers}
\subsubsection*{Provider GUID vs Group Provider GUID}
\subsubsection*{Checking correctness of logged events}
\subsubsection*{Automatization of event logging}
\subsection*{Service isolation}
\subsection*{Triggers}
\subsubsection*{searching for new triggers} YARA

\subsection*{Difference among configuration levels of telemtry}
\subsection*{Analysis of sent data over the channel to Microsfot backend services}



\section*{Results}
\section*{Conclusions}

%
%
%
%%%%% ABSTRACTS, AGRADECIMIENTOS Y DEDICATORIA
%\frontmatter
%\pagestyle{empty}
%\section*{Abstract (spanish version)}
%%\begin{center}
%\large \bf \runtitulo
%\end{center}
%\vspace{1cm}
\chapter*{\runtitulo}

\noindent La princesa Leia, líder del movimiento rebelde que desea reinstaurar la República en la galaxia en los tiempos ominosos del Imperio, es capturada por las malévolas Fuerzas Imperiales, capitaneadas por el implacable Darth Vader. El intrépido Luke Skywalker, ayudado por Han Solo, capitán de la nave espacial ``El Halcón Milenario'', y los androides, R2D2 y C3PO, serán los encargados de luchar contra el enemigo y rescatar a la princesa para volver a instaurar la justicia en el seno de la Galaxia (aprox. 200 palabras).

\bigskip

\noindent\textbf{Palabras claves:} Guerra, Rebelión, Wookie, Jedi, Fuerza, Imperio (no menos de 5).
%
%\cleardoublepage
%\section*{Abstract (english version)}
%Windows, one of the most popular OS, has a component called Telemetry. It collects information from the system with the goal of analyzing and fixing software \& hardware problems, improving the user experience, among others. The kind of information that can be obtained by this component is partially configurable in four different levels: security, basic, enhanced and full, being ”security” the level where less information is gathered and ”full” the opposite case.  

How Telemetry stores/process/administrates the information extracted? It employs a widely used framework called Event Tracer for Windows (ETW) [4]. Embedded not only in userland application but also in the kernel modules, the ETW framework has the goal of providing a common interface to log events and
therefore help to debug and log system operations, by instrumenting it.

In this work, we are going to analyze a part of the Windows kernel to better understand how Telemetry works from an internal perspective. This work will make windows analysts, IT admins or even windows users, more aware about the functionality of the Telemetry component helping to deal with privacy issues, bug fixing, knowledge of collected data, etc. Our analysis implies performing reverse engineering [5],[6] on the Telemetry component, which involves chal- lenging processes such as kernel debugging, dealing with undocumented kernel internal structures, reversing of bigger frameworks (i.e: ETW), binary libraries which lack symbols, etc. As part of the analysis we will also develop an in depth comparison between the differences among each level of Telemetry, stressing the contrast in the amount of events written, verbosity of information, etc. Finally, we will study how the communication between the Windows instance and the Microsoft backend servers is carried out.
 % OPCIONAL: comentar si no se quiere
%
%\cleardoublepage
%\center \  \section*{Greetings}
%\chapter*{Agradecimientos}

\noindent Lorem ipsum dolor sit amet, consectetur adipiscing elit. Fusce sapien ipsum, aliquet eget convallis at, adipiscing non odio. Donec porttitor tincidunt cursus. In tellus dui, varius sed scelerisque faucibus, sagittis non magna. Vestibulum ante ipsum primis in faucibus orci luctus et ultrices posuere cubilia Curae; Mauris et luctus justo. Class aptent taciti sociosqu ad litora torquent per conubia nostra, per inceptos himenaeos. Mauris sit amet purus massa, sed sodales justo. Mauris id mi sed orci porttitor dictum. Donec vitae mi non leo consectetur tempus vel et sapien. Curabitur enim quam, sollicitudin id iaculis id, congue euismod diam. Sed in eros nec urna lacinia porttitor ut vitae nulla. Ut mattis, erat et laoreet feugiat, lacus urna hendrerit nisi, at tincidunt dui justo at felis. Class aptent taciti sociosqu ad litora torquent per conubia nostra, per inceptos himenaeos. Ut iaculis euismod magna et consequat. Mauris eu augue in ipsum elementum dictum. Sed accumsan, velit vel vehicula dignissim, nibh tellus consequat metus, vel fringilla neque dolor in dolor. Aliquam ac justo ut lectus iaculis pharetra vitae sed turpis. Aliquam pulvinar lorem vel ipsum auctor et hendrerit nisl molestie. Donec id felis nec ante placerat vehicula. Sed lacus risus, aliquet vel facilisis eu, placerat vitae augue.
 % OPCIONAL: comentar si no se quiere
%
%\cleardoublepage
%\tableofcontents
%
%\mainmatter
%\pagestyle{headings}
%
%%%%% ACA VA EL CONTENIDO DE LA TESIS
%
%\chapter{La guerra de las galaxias}
%\section{Infancia y juventud}
%{\begin{small}%
%\begin{flushright}%
%\it
%There's nothing for me now.
%I want to learn the ways of\\ the Force and become a Jedi like my father. \\
%--Luke Skywalker
%\end{flushright}%
%\end{small}%
%\vspace{.5cm}}
%
%\section{Rescate de la princesa}
%{\begin{small}%
%\begin{flushright}%
%\it
%Here's where the fun begins!\\
%--Han Solo
%\end{flushright}%
%\end{small}%
%\vspace{.5cm}}
%
%\section{Sacrificio y victoria}
%{\begin{small}%
%\begin{flushright}%
%\it
%This will be a day long remembered.\\ It has seen the end of Kenobi.\\ It will soon see the end of the Rebellion.\\
%--Darth Vader
%\end{flushright}%
%\end{small}%
%\vspace{.5cm}}
%
%%% ...
%\chapter{El imperio contraataca}
%\chapter{El regreso del Jedi}
%
%%%%% BIBLIOGRAFIA
%\backmatter
%%\bibliography{tesis}

\end{document}
